\chapter{Research Timeline}

This chapter presents the planned schedule for the development and evaluation
of the hydrocarbon zone classification system, spanning five months from July
to November. The work is organized into eight phases that follow the
structure established in the methodology: literature consolidation, data
preparation, algorithmic development, interface construction, testing, and
final documentation.

\section{Phase Descriptions}

\subsection*{Literature Review \& Finalization (July)}
The opening phase consolidates the theoretical grounding of the research.
This includes a final review of the Pixler ratio framework \cite{ref13},
the Haworth--Whittaker composite indicators \cite{haworth1984, haworth1985},
and existing mud logging interpretation workflows. Any gaps identified in
the literature, particularly regarding the GOW, WBS, and GOR indicators are addressed before development begins. Deliverable: annotated bibliography and finalized Chapter 2.

\subsection*{Schema Definition (July--August)}
The structural architecture of the data is formally established by defining the expected column schema. This schema explicitly specifies the depth index alongside the seven required alkane components: Methane ($C_1$), Ethane ($C_2$), Propane ($C_3$), Iso-Butane ($iC_4$), Normal-Butane ($nC_4$), Iso-Pentane ($iC_5$), and Normal-Pentane ($nC_5$). Deliverable: Defined Schema Architecture

\subsection*{System Architecture Design (July--August)}
Before implementation begins, the three-layer architecture (parsing layer,
logic layer, UI layer) is designed in detail. This includes specifying the
data flow between layers, the callback structure for the Dash application,
and the classification threshold table for the majority-vote engine.
Deliverable: architecture document and pseudocode for the classification
pipeline.

\subsection*{Data Parsing \& Standardization Module (August)}
The first implementation phase covers the ingestion pipeline described in
Section~3.1: file upload handling, column validation, null filtering, and
z-score normalization.
Deliverable: working parsing module.

\subsection*{Algorithm \& Classification Logic (August--September)}
All fourteen hydrocarbon indicators (Pixler ratios, Haworth--Whittaker
ratios, TG, GOW, GOW No TG, WBS, GOR) are implemented and verified against
hand-computed reference values. The majority-vote classification engine is
then built on top of the indicator outputs, along with the tie-breaking
priority rule.
Deliverable: fully tested indicator and classification modules.

\subsection*{UI \& Visualization Development (September)}
The Dash application shell, upload modal, and results dashboard are built.
Area charts of each indicator trace and the colour-coded zone overlay are
implemented in Plotly. The interface is validated against the parsed dataset
from Phase 4.
Deliverable: integrated Dash application with end-to-end data flow.

\subsection*{System Testing \& Evaluation (October)}
The complete system is evaluated against the ground-truth well-test labels
using the classification metrics defined in Chapter~3: macro-averaged
Accuracy, Precision, Recall, and F1-Score. The confusion matrix per zone
class is computed, and performance on minority classes (e.g.\ oil-bearing
intervals) is examined in detail. Any misclassification patterns identified
are traced back to individual indicators for potential threshold adjustment.
Deliverable: evaluation results and performance tables for Chapter 4.

\subsection*{Writing, Revision \& Submission (October--November)}
Concurrent with testing in October, revisions are done in response
to supervisor feedback. Chapter 4 (Results) and Chapter 5 (Discussion and
Conclusion) are drafted once evaluation results are available. A full
manuscript review is conducted in early November, followed by final
proofreading and submission by the end of November.
Deliverable: submitted thesis manuscript.

% ------------------------------------------------------------------
\section{Gantt Chart}

Table~\ref{tab:gantt} presents the Gantt chart summarizing the timeline
described above. The horizontal axis spans twenty weeks divided into five
monthly blocks (July through November). Each bar represents the planned
active duration of its corresponding phase; overlapping bars indicate
concurrent work.

\bigskip

\begin{figure}[H]
\centering
\begin{ganttchart}[
    hgrid,
    vgrid,
    x unit=2.2cm,
    y unit title=0.65cm,
    y unit chart=0.7cm,
    title height=1,
    bar height=0.5,
    bar/.append style={fill=blue!25, rounded corners=2pt},
    title label font=\small\bfseries,
    bar label font=\small,
    bar label node/.append style={text width=5.2cm, align=left},
]{1}{5}
  \gantttitle{July}{1}
  \gantttitle{August}{1}
  \gantttitle{September}{1}
  \gantttitle{October}{1}
  \gantttitle{November}{1} \\
  \ganttbar{1.\ Literature Review}{1}{1} \\
  \ganttbar{2.\ Schema Definition}{1}{2} \\
  \ganttbar{3.\ Architecture Design}{1}{2} \\
  \ganttbar{4.\ Parsing Module}{2}{2} \\
  \ganttbar{5.\ Algorithm \& Classification}{2}{3} \\
  \ganttbar{6.\ UI \& Visualization}{3}{3} \\
  \ganttbar{7.\ Testing \& Evaluation}{4}{4} \\
  \ganttbar{8.\ Writing \& Submission}{4}{5}
\end{ganttchart}
\caption{Planned research timeline, July--November 2025.}
\label{tab:gantt}
\end{figure}